% TOP_PROBLEM: {"id": 30007280, "title": "Rationality of the canonical unramified Arthur filtration", "category_id": 1, "category": {"id": 1, "name": "number_theory", "display_name": "Number Theory", "description": "Properties of integers, prime numbers, Diophantine equations.", "slug": "number-theory", "order_index": 1, "created_at": "2026-07-31T15:26:25.670Z"}, "status": "open"}
% FIRST_POSTED: 2026-10-06
% ATTEMPT_STATUS: solved
% OPENAI_CLAIM: {"schema_version":1,"source_repo":"https://github.com/openai/math","source_commit":"adc7f1241b42e322a6451854ab7e4b4c146bf78a","release_date":"2026-10-06","match":"full","resolution":"proved","independently_reviewed":false,"second_pass":"agreed","families":[{"family":"014","title":"Restricted geometric Langlands, global Arthur enhancements, and generic Ramanujan","subject":"Number theory","manuscripts":[{"title":"Rationality of the Canonical Unramified Arthur Filtration","pdf_path":"preprints/Rationality-of-the-Canonical-Unramified-Arthur-Filtration-September-24-2026/paper.pdf","readme_path":"preprints/Rationality-of-the-Canonical-Unramified-Arthur-Filtration-September-24-2026/README.md","date":"2026-09-24","theorem_ref":"intro:main, intro:main-map (sections/01-introduction.tex:118)","statement_sources":[{"path":"preprints/Rationality-of-the-Canonical-Unramified-Arthur-Filtration-September-24-2026/build/sections/01-introduction.tex","line":118,"label":"intro:main, intro:main-map"}]}],"lean":null,"reasoning_trace":null}]}
\documentclass[11pt]{article}
\usepackage[margin=25mm]{geometry}
\usepackage{fontspec}
\setmainfont{Latin Modern Roman}
\usepackage{amsmath,amssymb,mathtools}
\usepackage{unicode-math}
\setmathfont{Latin Modern Math}
\usepackage{enumitem}
\usepackage{xurl,hyperref,bookmark}
\hypersetup{colorlinks=true,urlcolor=blue}
\setcounter{secnumdepth}{0}
\setlength{\parindent}{0pt}
\setlength{\parskip}{5pt}
\setlength{\emergencystretch}{3em}
\begin{document}
\section*{Rationality of the canonical unramified Arthur filtration}

\subsection{OpenAI's claim}
OpenAI presents the following claims relevant to this question, with manuscript numbers following paper order in Sources. In manuscript 1 (family 014; 2026-09-24), the asserted result is that the canonical unramified Arthur filtration on all finitely supported automorphic functions is rational under the restricted theory's characteristic hypotheses for split connected semisimple groups.

\subsection{Mathematical statement of the claimed result}
\paragraph{Manuscript 1.}
Let $X/\mathbb F_q$ be smooth projective geometrically connected and $G$ split connected semisimple. On $H=G_{\overline{\mathbb F}_q}$ assume an invariant nondegenerate Lie-algebra form nondegenerate on all Levi centers, Chevalley restriction for every Levi, semisimple Lie-algebra centralizers being Levi, and every nilpotent element over any extension lying in a parabolic unipotent-radical Lie algebra. For $E=\overline{\mathbb Q}_{\ell}$, $\ell\ne\operatorname{char}\mathbb F_q$, let $S$ be bundle isomorphism classes over $\mathbb F_q$ and $\mathcal A_E=E^{(S)}$. For a closed invariant $Y$ in the pinned dual-group nilpotent cone, use the derived-Satake support category $\mathcal C_Y$ and Frobenius $\Phi$ to define
\[
 \mathcal F_Y\mathcal A_E=\operatorname{im}(H^0\operatorname{Tr}(\Phi,\mathcal C_Y)\to\mathcal A_E).
\]
The claim says this entire canonical filtration is rational:
\[
 E\otimes_{\mathbb Q}(\mathcal A_{\mathbb Q}\cap\mathcal F_Y\mathcal A_E)
 \xrightarrow{\sim}\mathcal F_Y\mathcal A_E
\]
for every such $Y$. The intersection is inside $\mathcal A_E$. Noncuspidal functions are included, and supported traces are defined by their image without assuming concentration in degree zero.

\subsection{Outline of the argument}
For manuscript 1 (family 014; 2026-09-24), a bilinear detector tests supported spectral local cohomology against finite evaluations. Orbit distributions, compact Weil lifts, and noncancellation separate the nilpotent supports, and rational descent then includes the noncuspidal functions.

\subsection{Scope relative to this problem}
The result supplies the conclusion in this problem with its displayed hypotheses and quantifiers. The broader statements and applications in the manuscript are kept separate from the particular question above. The supporting result is manuscript 1 (family 014; 2026-09-24).

\subsection{Formal verification}
No Lean formalization is listed at this snapshot for manuscript 1 (family 014; 2026-09-24).

This site has not executed the Lean code. Described statement coverage is not an independent proof check.

\subsection{Status caveat}
The manuscripts were produced by an internal OpenAI model. OpenAI cautions that results without formalization can contain mistakes. This site has not independently checked the proofs; the label records the scope of OpenAI's claim.

\subsection{Sources}
\begin{itemize}
\item OpenAI, \emph{Rationality of the Canonical Unramified Arthur Filtration}, 2026-09-24; citation key \texttt{OAI:Rationality-of-the-Canonical-Unramified-Arthur-Filtration-September-24-2026}. \href{https://github.com/openai/math/blob/adc7f1241b42e322a6451854ab7e4b4c146bf78a/preprints/Rationality-of-the-Canonical-Unramified-Arthur-Filtration-September-24-2026/paper.pdf}{PDF (GitHub)}; \href{https://github.com/openai/math/raw/adc7f1241b42e322a6451854ab7e4b4c146bf78a/preprints/Rationality-of-the-Canonical-Unramified-Arthur-Filtration-September-24-2026/paper.pdf}{PDF (direct)}; \href{https://github.com/openai/math/blob/adc7f1241b42e322a6451854ab7e4b4c146bf78a/preprints/Rationality-of-the-Canonical-Unramified-Arthur-Filtration-September-24-2026/README.md}{Manuscript page and citation}; \href{https://github.com/openai/math/blob/main/preprints/Rationality-of-the-Canonical-Unramified-Arthur-Filtration-September-24-2026/paper.pdf}{Latest version}; Statement source: \href{https://github.com/openai/math/blob/adc7f1241b42e322a6451854ab7e4b4c146bf78a/preprints/Rationality-of-the-Canonical-Unramified-Arthur-Filtration-September-24-2026/build/sections/01-introduction.tex\#L118}{intro:main, intro:main-map}.
\item \href{https://github.com/openai/math/blob/adc7f1241b42e322a6451854ab7e4b4c146bf78a/CONTENTS.md}{OpenAI catalogue entry}.
\end{itemize}
\end{document}
